\section{Release Scope, Capability, and Directory Map}

VENUS is one release with two implementation trees:
\path{VENUS/cpu} and \path{VENUS/gpu}.  ``CPU'' below means the broad
reference implementation, including serial, \code{std::thread}, OpenMP, and
the supported MPI phase space decomposition.  ``GPU'' means the CUDA/MPI WPIC
implementation.  The GPU tree accelerates the weighted particle--mesh family, including its
material and subhalo collision paths. Grid Einstein--Vlasov, grid
semi-Lagrangian Vlasov--Poisson, canonical-orbit Boltzmann--Poisson, and
SphDSMC use the CPU implementation. Choose the physical task in the capability
table before selecting its CPU or GPU executable.

The manual is organized by physical task.  Each module entry identifies its
equations, numerical representation, task-specific source and executable,
applicable parameters, IC/restart and output contract, parallel model,
visualization path, and scientific limitations.  Shared kernels are numerical functions called by several tasks. Their
algorithms are explained where introduced, and the task executable supplies
the complete run loop that calls them.  A calculation intended for quantitative interpretation still
requires the module-specific convergence, conservation, boundary, locality,
and parallel-equivalence checks stated in its chapter.

\subsection{Capability map}

Every path in the next table is relative to the repository's
\path{VENUS/cpu} directory.  For example,
\path{include/ev/config.hpp} means the repository path
\path{cpu/include/ev/config.hpp}; this convention is confined to this one
directory table.

\begin{center}
\small
\begin{tabularx}{\linewidth}{@{}>{\raggedright\arraybackslash}p{0.28\linewidth}XX@{}}
\toprule
Physical/numerical module & CPU implementation & GPU implementation \\
\midrule
Einstein--Vlasov & grid semi-Lagrangian modular and stand-alone solvers;
separate Schwarzschild/ingoing-EF weighted-marker tasks; bound static-shell IC
and same-event coordinate mapper & not implemented \\
Static/expanding Vlasov--Poisson & nodal, finite-volume, conservative SLDG,
Hamiltonian/canonical, and WPIC paths & WPIC kick--drift--kick with CUDA
deposition, force interpolation, and optional MPI \\
Boltzmann--Poisson & grid collision operators, canonical events, WPIC, and
SphDSMC & WPIC material and subhalo collision operators \\
Constant SIDM & grid/FP or sampled finite-angle paths as configured; WPIC and
SphDSMC & material-pair Monte Carlo \\
Velocity-dependent SIDM & power-law and Yukawa/Born models in their documented
CPU paths & analytic softened power-law and Yukawa--Born material kernels with
GPU pair sampling; interpolation tables belong to the separate subhalo model \\
Subhaloes & unequal-mass Plummer scattering/capture in WPIC & unequal-mass
elastic/capture events with local pair majorants \\
Collision timestep & mesh-wide reference mode; independent CPU subhalo cells;
independent equal-mass constant-SIDM shells & mesh-wide reference mode;
independent subhalo cells; independent material shells for every GPU
cross-section law \\
Parallelism & \code{std::thread}, OpenMP; MPI for supported grid/transport
paths; collision guards are fail-closed where unsupported & CUDA plus optional
MPI domain/rank work; multi-GPU checks require a scheduler allocation \\
\bottomrule
\end{tabularx}
\end{center}

\subsection{Directory map}

\codesource{Release directory topology}{\path{Makefile}, \path{cpu/Makefile}, and \path{gpu/CMakeLists.txt}}{Maps the source, build, documentation and output roots maintained by the release}
\begin{lstlisting}
VENUS/
  cpu/       CPU source, task front ends, examples and tools
  gpu/       CUDA/MPI WPIC source, examples, scripts and documentation
  docs/      compiled manual, documentation index and module guides
  notes/     editable manual sources and solver/script catalogues
  scripts/   release-wide utility and verification scripts
  output/    intentionally empty CPU/GPU run roots
\end{lstlisting}

\paragraph{CPU subtree.}
\codesource{CPU source-tree map}{\path{cpu/Makefile}}{Identifies the CPU task fronts, shared drivers and headers, tools, examples, scheduler templates, and output root owned by the CPU build}
\begin{lstlisting}
cpu/
  src/apps/             one main for each Eulerian/BP/SphDSMC/WPIC task
  src/drivers/          shared WPIC application engines
  include/ev/           shared physics, algorithms, parallelism, I/O
  tools/                IC, restart, analysis, plotting, and GIF utilities
  python/ev_solver/     readable Python reference implementation
  examples/initial_data/
  scripts/slurm/        portable scheduler templates
  docs/                 detailed CPU method and cluster notes
  ../output/cpu/        root-level CPU run products
\end{lstlisting}

\paragraph{GPU subtree.}
\codesource{GPU source-tree map}{\path{gpu/CMakeLists.txt}}{Identifies the CUDA interfaces, kernels, task fronts, benchmarks, examples, visualization helper, and output root registered by the GPU build}
\begin{lstlisting}
gpu/
  include/venus_gpu/    public WPIC/collision/restart/MPI interfaces
  src/                  CUDA kernels and the shared task engine
  tools/                one main per GPU task and restart/table utilities
  benchmarks/           focused VP and material-SIDM timing programs
  scripts/              GPU plot/GIF helper
  docs/                 GPU equations, algorithms, API, and MPI notes
  ../output/gpu/        root-level GPU run products
\end{lstlisting}

\paragraph{Build and first-run map.}
\codesource{Release build and first-run entry points}{\path{Makefile}, \path{cpu/Makefile}, \path{gpu/CMakeLists.txt}, \path{cpu/src/apps/vlasov_poisson/eulerian.cpp}, and \path{gpu/tools/venus_gpu_vlasov_poisson.cu}}{Builds the supported CPU and GPU task fronts and demonstrates their fail-closed Vlasov--Poisson launch interfaces}
\begin{lstlisting}[language=bash]
make -C cpu organized-solvers
make -C cpu cpu-task-solvers-openmp
make -C cpu cpu-task-solvers-mpi       # Eulerian EV/VP only
make gpu GPU_ARCH=80
make gpu-mpi GPU_ARCH=80

cpu/venus_cpu_vlasov_poisson_eulerian \
  --config cpu/examples/modular/vp_expanding_nfw_gamma2_collisionless.ini
build/gpu/venus_gpu_vlasov_poisson --help
\end{lstlisting}
The detailed installation and launch sections below explain prerequisites,
architecture selection, configuration, output roots, and scheduler resources.

New calculations should use the task-specific interfaces catalogued below.
Compatibility dispatchers remain only so older inputs can be interpreted.

\section{Scope and Directory Map}

VENUS evolves a spherically symmetric distribution function
\[
  f=f(t,r,w,L),
\]
where \(r\) is radius, \(w\) is the radial momentum or peculiar radial velocity
depending on the model, and \(L=|\bm r\times\bm v|\) is angular momentum
magnitude.
\ifvenusfullcpu
The CPU tree provides task-specific C++ front ends:
\begin{itemize}
\item \code{venus_cpu_einstein_vlasov_eulerian} and
\code{venus_cpu_vlasov_poisson_eulerian} lock the grid equation;
\item the \code{venus_cpu_boltzmann_poisson_*} names lock the selected
Fokker--Planck/Landau or finite-angle backend;
\item \code{venus_cpu_sidm_sphdsmc_*} separates constant and
velocity-dependent persistent-marker SIDM; and
\item the six \code{venus_cpu_*_wpic*} names separate relativistic,
collisionless, material-SIDM, and subhalo marker tasks.
\end{itemize}
The older \code{venus_modular_solver*}, \code{venus_einstein_vlasov*}, and
\code{run_*} programs are compatibility or specialist interfaces.
\else
The marker paths documented here use task-specific names:
\begin{itemize}
\item SphDSMC uses \code{venus_cpu_sidm_sphdsmc_constant} or
\code{venus_cpu_sidm_sphdsmc_velocity_dependent}.
\item WPIC uses the physical-task names in
Section~\ref{sec:task-solver-catalogue}.
\end{itemize}
The numerical descriptions and limitations appear in the CPU chapters below.
\fi

\begin{center}
\footnotesize
\begin{tabular}{@{}>{\raggedright\arraybackslash}p{0.38\linewidth}>{\raggedright\arraybackslash}p{0.54\linewidth}@{}}
\toprule
Path & Contents \\
\midrule
\path{src/apps/} & separate task mains for Eulerian, BP, SphDSMC, and WPIC applications \\
\path{src/drivers/} & shared WPIC application engines used by the task mains \\
\ifvenusfullcpu
\path{src/venus_einstein_vlasov.cpp} & self-contained EV cluster source \\
\fi
\path{include/ev/common.hpp} & runtime parameter structure and model selectors \\
\path{include/ev/config.hpp} & config, CLI, environment, and launch-safety parser \\
\path{include/ev/grids.hpp} & grids, indexing, control-volume edges, quadrature \\
\ifvenusfullcpu
\path{include/ev/moments.hpp}, \path{include/ev/metric.hpp} & EV moments and metric constraints \\
\path{include/ev/vlasov_poisson.hpp} & static/expanding VP field solve, transport, finite-volume option, inflow \\
\path{include/ev/conservative_sldg_transport.hpp} & conservative source-cell VP SLDG map, polygon deposition, midpoint field, and delayed carrier \\
\path{include/ev/hamiltonian_sldg_transport.hpp}, \path{include/ev/hamiltonian_orbit_reference.hpp} & Hamiltonian-aligned reconstruction and positive orbit-average reference \\
\path{include/ev/boltzmann_poisson.hpp} & collision-module dispatch \\
\path{include/ev/finite_angle_boltzmann.hpp} & shared cross-section and collision-kinematic utilities \\
\path{include/ev/canonical_orbit_phase_transport.hpp} & finite-volume orbit atlas, frozen-field phase transport, collision lifts, and action-preserving field rebuilds \\
\path{include/ev/canonical_boltzmann_poisson_event.hpp} & adaptive field/collision event schedule and conservation budgets \\
\path{include/ev/canonical_particle_ntc_collision.hpp}, \path{include/ev/canonical_reversible_collision.hpp} & parcel-owned NTC and increment-only reversible collision coupling \\
\path{include/ev/canonical_open_carrier_compression.hpp} & positive local quadrature reduction for continuous collision packets \\
\path{include/ev/canonical_boltzmann_poisson_restart.hpp} & canonical atlas, packet, random-stream, and event-clock restarts \\
\else
\path{include/ev/weighted_vlasov_pic.hpp} & weighted marker initialization, radial deposition, orbit integration, delayed inner return, and restarts \\
\path{include/ev/weighted_vlasov_pic_collisions.hpp} & local NTC scattering, mean-free-path shell refinement, and sampled collision-time estimates \\
\path{include/ev/weighted_vlasov_pic_subhalo_collisions.hpp} & unequal-mass whole-object Plummer scattering/capture, remnant merging, and internal-energy ledger \\
\path{include/ev/weighted_vlasov_pic_event_scheduler.hpp} & independent field and collision clocks \\
\path{include/ev/weighted_vlasov_pic_phase_remap.hpp} & conservative marker-to-grid reconstruction \\
\path{include/ev/weighted_pic_boltzmann_event.hpp} & integrated frozen-field transport, NTC collision, and reconstruction events \\
\path{include/ev/spherical_dsmc_collision.hpp} & direct finite-angle scattering of continuous marker velocities \\
\path{include/ev/spherical_dsmc_solver.hpp} & marker gravity, radial collision shells, stiff-orbit integration, budgets, and restarts \\
\fi
\ifvenusfullcpu
\path{include/ev/advection.hpp} and \path{include/ev/timestep.hpp} & characteristic pullback, boundaries, adaptive \(w\), timestep scans \\
\fi
\path{include/ev/mpi.hpp} and \path{include/ev/parallel.hpp} & MPI \(L\)-slab decomposition and OpenMP/std::thread wrappers \\
\path{include/ev/output.hpp} & CSV output, metadata, phase projections, restarts \\
\path{examples/} & publishable example parameter files and binary ICs \\
\path{scripts/slurm/} & IAS and Indian Institute of Technology Kanpur (IITK)
Slurm templates \\
\path{python/ev_solver/} & Python reference solver, reader, plotting, and IC helpers \\
\path{tools/} & plotting, GIF, IC, and diagnostic scripts \\
\bottomrule
\end{tabular}
\end{center}

The modular physics is implemented in \path{include/ev/*.hpp}. Follow the
selected task main into these headers to connect its input choices to the
field solve, transport, collisions, and output operations described here.

\section{Build and Quick Start}

Build from the \path{VENUS/cpu} directory:
\codesource{CPU task-front build}{\path{cpu/Makefile}}{Builds the portable, OpenMP, MPI-enabled collisionless, and debug CPU executable sets documented below}
\begin{lstlisting}
cd VENUS/cpu
make organized-solvers        # all portable task-specific fronts
make cpu-task-solvers-openmp  # grid/SphDSMC OpenMP fronts
make cpu-task-solvers-mpi     # Eulerian EV/VP MPI+OpenMP fronts
make debug       # unoptimized debug builds
\end{lstlisting}
The build targets produce the following executables:
\begin{itemize}
\ifvenusfullcpu
\item portable: the unsuffixed \code{venus_cpu_*} task names;
\item OpenMP: the grid/SphDSMC task names ending in \code{_openmp};
\item MPI+OpenMP: \path{venus_cpu_einstein_vlasov_eulerian_mpi} and
\path{venus_cpu_vlasov_poisson_eulerian_mpi}.
\else
\item SphDSMC: the \code{venus_cpu_sidm_sphdsmc_*} task names;
\item WPIC: the six task-specific names in the catalogue.
\fi
\end{itemize}

Useful compiler overrides are:
\codesource{CPU compiler and parallel-runtime overrides}{\path{cpu/Makefile}}{Selects the C++ or MPI compiler wrapper and its matching OpenMP flag without changing task ownership}
\begin{lstlisting}
make cpu-task-solvers-openmp CXX=g++ OPENMP_FLAGS=-fopenmp
make cpu-task-solvers-mpi MPICXX=mpicxx OPENMP_FLAGS=-fopenmp
\end{lstlisting}
On macOS with Apple Clang, \code{make organized-solvers} supports systems
whose compiler lacks OpenMP.

Run the lightweight startup and I/O checks with:
\codesource{Packaged CPU setup check}{\path{cpu/Makefile}}{Builds the compatibility solver and runs its maintained zero-time startup and output smoke checks}
\begin{lstlisting}
make check-setup
\end{lstlisting}
\ifvenusfullcpu
\codesource{Eulerian Vlasov--Poisson zero-time check}{\path{cpu/src/apps/vlasov_poisson/eulerian.cpp}, \path{cpu/include/ev/solver.hpp}, and \path{cpu/examples/modular/vp_bound_infalling_gaussian_shell_train.ini}}{Runs the fail-closed collisionless Newtonian frontend through initialization and output without taking an evolution step}
\begin{lstlisting}
./venus_cpu_vlasov_poisson_eulerian \
  --config examples/modular/vp_bound_infalling_gaussian_shell_train.ini \
  --t-end 0 --output-dir output/check_vp_shell --threads 2
\end{lstlisting}
\else
\codesource{Eulerian Fokker--Planck zero-time check}{\path{cpu/src/apps/boltzmann_poisson/fokker_planck.cpp}, \path{cpu/include/ev/solver.hpp}, and \path{cpu/examples/modular/bp_sidm_corecollapse_yukawa.ini}}{Runs the fail-closed Fokker--Planck Boltzmann--Poisson frontend through initialization and output without taking an evolution step}
\begin{lstlisting}
./venus_cpu_boltzmann_poisson_fokker_planck \
  --config examples/modular/bp_sidm_corecollapse_yukawa.ini \
  --t-end 0 --output-dir output/check_bp_sidm --threads 2
\end{lstlisting}
\fi

\ifvenusfullcpu
Run the supplied examples with:
\codesource{Packaged Eulerian EV, VP, and Fokker--Planck examples}{\path{cpu/src/apps/einstein_vlasov/eulerian.cpp}, \path{cpu/src/apps/vlasov_poisson/eulerian.cpp}, \path{cpu/src/apps/boltzmann_poisson/fokker_planck.cpp}, and \path{cpu/examples}}{Launches the task-locked CPU frontends with the maintained relativistic, collisionless Newtonian, and collisional Newtonian example configurations}
\begin{lstlisting}
./venus_cpu_einstein_vlasov_eulerian \
  --config examples/einstein_vlasov_high_compactness_horizon_formation.ini

./venus_cpu_einstein_vlasov_eulerian \
  --config examples/einstein_vlasov_low_compactness_virialized_halo.ini

./venus_cpu_vlasov_poisson_eulerian \
  --config examples/modular/vp_expanding_nfw_gamma2_collisionless.ini

./venus_cpu_boltzmann_poisson_fokker_planck \
  --config examples/modular/bp_expanding_nfw_gamma1p5_weak_collision.ini

./venus_cpu_boltzmann_poisson_fokker_planck \
  --config examples/modular/bp_sidm_corecollapse_yukawa.ini

./venus_cpu_vlasov_poisson_eulerian \
  --config examples/modular/vp_bound_infalling_gaussian_shell_train.ini
\end{lstlisting}

The solver prints the resolved model and parallel layout at startup.  For
example,
\codesource{Resolved solver-layout startup banner}{\path{cpu/include/ev/solver.hpp}}{Reports the selected physics model, parallel backend, threads per rank, MPI size, and global angular momentum grid size after setup}
\begin{lstlisting}
Solver model=vlasov_poisson parallel backend=openmp threads_per_rank=8 mpi_size=4 global_NL=72
\end{lstlisting}
is a healthy hybrid EV/VP launch.  For BP collisions, \code{mpi_size} must be
one in this release.
\else
The module chapters below begin with the required plain-language physics
flowchart and then give a complete runnable command plus the corresponding
\path{params.ini} or runner parameters.  For a full SphDSMC example, use the
command above without the \code{--t-end 0} startup-check override.  For WPIC, build the
runner named in its chapter and use the shown Frank-core or expanding
    power-law command.  Both marker modules use one process with threaded
marker/orbit work; do not launch multiple MPI ranks for their collision steps.
\fi

\section{Configuration Files}

Input files group key--value entries into named sections for readability.
The parser reads the keys themselves, so section headings organize the file
while leaving each key's meaning unchanged.  Both \code{key = value} and
\code{key: value} are accepted.  Unquoted \# and semicolon comments are
stripped.  Underscores and hyphens are equivalent, so
\code{bp_collision_flux}, \code{bp-collision-flux}, and the command-line
option \code{--bp-collision-flux} all refer to the same parameter.

The precedence order is:
\begin{enumerate}
\item compiled defaults in \path{include/ev/common.hpp};
\item the selected config file;
\item environment overrides: \code{EV_CONFIG}, \code{EV_OUTPUT_DIR},
\code{EV_SOLVER_OUTPUT_DIR}, \code{EV_RESTART_DIR}, \code{EV_RESTART_FROM},
\code{EV_THREADS}, \code{EV_SOLVER_THREADS}, and
\code{EV_ALLOW_MULTI_RANK_LAUNCH};
\item command-line overrides.
\end{enumerate}
In the absence of \code{--config}, the solver searches for
\path{input/params.ini} and then \path{params.ini} in the current directory and
next to the executable.  Relative paths inside a config file are resolved
relative to that config file.

The standard sections are:
\begin{itemize}
\item \code{[physics]}: \code{physics_model}, VP background, Poisson source,
and interpolation mode.
\item \code{[collisions]}: BP collision operator, cross section, bath,
coefficient refresh, flux, and integrator.
\item \code{[initial_data]}: analytic or EVICRWL1 binary initial data, optional
tracer, and mass normalization.
\item \code{[grid_for_analytic_initial_data]} or
\code{[grid_for_binary_initial_data]}: resolution and domain.
\item \code{[adaptive_w]}: optional one-way expansion of the represented
momentum domain.  In VENUS this is disabled for binary IC workflows.
\item \code{[time_integration]}: integrator, CFL, final time, and step cap.
\item \code{[boundaries]}: inner, outer, momentum, and central-mass controls.
\item \code{[metric_and_stopping]}: metric floors, compactness stop, active
cell floors, and distribution floor.
\item \code{[output]}: output and restart directories and write cadence.
\item \code{[parallelism]}: thread count and accidental multi-rank guard.
\end{itemize}
